\chapter{Tecnologías y Arquitectura de Red}

Una vez establecido el posicionamiento del proyecto respecto del estado del arte, corresponde describir las decisiones tecnológicas concretas que sostienen la implementación del prototipo, así como la arquitectura de red que enmarca su funcionamiento. A diferencia de un desarrollo de software genérico, este proyecto está condicionado por un factor externo determinante: la aplicación no se ejecuta de manera autónoma, sino embebida dentro del panel de administración de Tiendanube, lo que impone restricciones técnicas concretas sobre el lenguaje del frontend y sobre el modo en que la aplicación debe ser servida durante el desarrollo.

\section{Criterios de selección del stack tecnológico}
\label{sec:criterios-seleccion}

Las decisiones tecnológicas del proyecto responden a dos tipos de criterios complementarios. El primero es una restricción impuesta por la plataforma de integración: Tiendanube exige que las aplicaciones integradas a su panel de administración estén construidas sobre React, dado que tanto su \emph{Software Development Kit} (SDK) de comunicación con el panel, \texttt{@tiendanube/nexo} (en adelante, Nexo), como su sistema de diseño oficial, Nimbus, fueron desarrollados específicamente para ese framework. Esta condición no es una preferencia del equipo, sino un requisito no funcional externo que determina de manera directa la tecnología del frontend.

El segundo criterio es una decisión propia del equipo: dado que el frontend debía ser necesariamente JavaScript, se optó por utilizar Node.js también en el backend, de modo de mantener un único lenguaje de programación a lo largo de todo el proyecto. Esta decisión reduce la carga cognitiva de un equipo reducido de dos personas, que debe alternar constantemente entre ambas capas de la aplicación, y evita la duplicación de herramientas de tooling (linters, formateadores, gestores de paquetes) que implicaría sostener dos ecosistemas de lenguajes distintos.

A estos dos criterios se suma una consideración de alcance ya establecida: el proyecto no contempla un despliegue productivo a escala industrial, sino la construcción de un prototipo funcional. Esta definición, presentada en el Alcance del proyecto, habilita decisiones de infraestructura más simples que las que se justificarían en un producto en producción, como se desarrolla en las secciones siguientes.

\section{Organización del código: monorepo}

El proyecto se estructura como un \emph{monorepo}: backend y frontend conviven en un único repositorio, organizados en las carpetas \texttt{backend/} y \texttt{frontend/} respectivamente. Esta decisión responde al fuerte acoplamiento funcional entre ambas capas. Dado que gran parte de los cambios en el proyecto involucra simultáneamente la incorporación o modificación de un endpoint en el backend y su consumo correspondiente en el frontend, un esquema de monorepo permite que ambos cambios se integren en un único commit, evitando que las dos partes queden desincronizadas entre sí, un riesgo real en un esquema de múltiplos repositorios cuando cada uno versiona de forma independiente.

La alternativa de un esquema multi-repo, con versionado y ciclos de despliegue independientes para cada capa, resulta apropiada cuando ambas partes deben evolucionar y desplegarse de forma desacoplada, típicamente en el marco de un despliegue productivo con múltiples equipos. Dado que, como se mencionó, el proyecto no contempla un despliegue productivo a escala industrial sino un prototipo funcional, esa ventaja del multi-repo no aplica en este contexto, mientras que sí se mantiene el costo de coordinación que impone. En consecuencia, el monorepo resulta la opción que mejor equilibra simplicidad operativa y consistencia entre las dos capas de la aplicación.

\section{Stack tecnológico}

\subsection{Frontend: React, Vite y el ecosistema de Tiendanube}

El frontend está desarrollado en React (versión \texttt{19.2.7}), servido mediante Vite (versión \texttt{8.1.1}) como herramienta de \emph{build} y servidor de desarrollo. Como se indicó, la elección de React no es opcional: Tiendanube exige este framework para las aplicaciones integradas a su panel de administración, dado que su SDK Nexo y su sistema de diseño Nimbus están construidos específicamente para él. Nexo es responsable de establecer la comunicación entre la aplicación, ejecutada dentro de un \emph{iframe}, y la ventana del panel de administración que la contiene, mientras que Nimbus provee los componentes visuales con los que la interfaz del prototipo mantiene consistencia estética con el resto del panel de Tiendanube.

Sobre esa base obligatoria de React, Vite fue seleccionado como herramienta de desarrollo por sus características propias: a diferencia de herramientas de \emph{build} más antiguas, Vite sirve el código fuente mediante módulos ES nativos del navegador durante el desarrollo, lo que se traduce en un servidor de arranque prácticamente instantáneo y en actualización en caliente (\emph{Hot Module Replacement}) de los cambios, agilizando el ciclo de desarrollo iterativo propio de un prototipo.

La integración entre Vite, Nexo y el entorno de túneles necesario para el desarrollo (descripto en la Sección~\ref{sec:arquitectura-red}) no resultó automática. Nexo fue diseñado originalmente para entornos Node.js y referencia en su código la variable global \texttt{global}, inexistente en el navegador; sin una configuración adicional, la aplicación fallaba al iniciar con el error \texttt{Uncaught ReferenceError: global is not defined}, que se traducía en una pantalla en blanco dentro del \emph{iframe}. La solución consistió en instruir a Vite, mediante la opción \texttt{define: \{ global: 'globalThis' \}}, para reemplazar en el código empaquetado todas las referencias a \texttt{global} por \texttt{globalThis}, sí disponible en los navegadores modernos. De manera similar, Vite bloquea por defecto las peticiones entrantes desde hosts externos al del propio servidor de desarrollo, lo cual entra en conflicto directo con el uso de un túnel público (Sección~\ref{sec:arquitectura-red}): al acceder a la aplicación mediante su URL pública, el navegador recibía el error \texttt{Blocked request. This host is not allowed}. Esto se resolvió habilitando explícitamente el dominio del proveedor de túnel en \texttt{server.allowedHosts}. Ambos ajustes ilustran que, más allá de la elección del framework y la herramienta de \emph{build}, buena parte del trabajo técnico consistió en conciliar un SDK pensado para Node.js con las restricciones propias del navegador y de la arquitectura de red descripta más adelante.

\subsection{Backend: Node.js y Express}
\label{sec:backend-node-express}

El backend está desarrollado en Node.js, en su versión 18 o superior, único requisito de versión documentado para el proyecto, utilizando Express como framework para la construcción de la API REST. Como se explicó en la Sección~\ref{sec:criterios-seleccion}, la elección de Node.js responde directamente a la necesidad de mantener un único lenguaje de programación en todo el proyecto, una vez que el frontend quedó determinado por el requisito de React de Tiendanube. Express fue seleccionado por tratarse de un framework minimalista y no dogmático, ampliamente adoptado para la construcción de API REST, que permite definir rutas y middlewares con una curva de aprendizaje reducida, adecuada al alcance de un prototipo desarrollado por un equipo de dos personas.

El backend cumple dos funciones centrales. Por un lado, actúa como intermediario autenticado entre el frontend y la API de Tiendanube: es quien retiene las credenciales de la aplicación (\texttt{Client ID} y \texttt{Client Secret}) y ejecuta el intercambio OAuth 2.0 descripto en la Sección~\ref{sec:estrategia-datos-ref}, de modo que dichas credenciales nunca se exponen al frontend ni al \emph{iframe} embebido en el navegador del usuario. Por otro lado, expone los \emph{webhooks} de privacidad que Tiendanube exige a toda aplicación integrada a su ecosistema (\texttt{/api/webhooks/store-redact}, \texttt{customers-redact} y \texttt{customers-data-request}), mediante los cuales la plataforma notifica solicitudes de eliminación de datos en cumplimiento de normativas de protección de datos personales.

\subsection{Persistencia: PostgreSQL}

Para la capa de persistencia se seleccionó PostgreSQL como motor de base de datos relacional. Al momento de esta entrega, la base de datos aún no fue incorporada al proyecto; el prototipo opera exclusivamente contra la API de Tiendanube y datos en memoria, y la persistencia se incorporará en una etapa posterior del desarrollo. La elección de un motor relacional responde a la naturaleza de las entidades centrales del dominio ---pedidos, productos, variantes y clientes---, que presentan relaciones estructuradas y bien definidas entre sí, tal como se describió al detallar los recursos consumidos de la API de Tiendanube en el Marco Teórico. El diseño del modelo de datos correspondiente excede el alcance de la presente sección y será abordado en el capítulo dedicado a dicho modelo.

\subsection{Túnel de desarrollo: ngrok}

El uso de ngrok como herramienta de túnel HTTPS es una consecuencia directa de la arquitectura de integración con Tiendanube, y se desarrolla en detalle en la Sección~\ref{sec:arquitectura-red}. Cabe adelantar que su necesidad no es una elección arbitraria de infraestructura, sino la solución a una restricción concreta: Tiendanube no admite que una aplicación de prueba se sirva desde una URL \texttt{localhost}, y requiere en su lugar una URL pública con HTTPS. ngrok resuelve esta necesidad generando, a partir del servidor de desarrollo local, una URL pública que actúa como túnel hacia dicho servidor, sin requerir desplegar la aplicación en un servidor propio durante la etapa de desarrollo.

Esta solución es adecuada al carácter de prototipo del proyecto, dado que evita construir infraestructura de hosting propia mientras el despliegue productivo permanece fuera de alcance. Tiene, no obstante, una limitación relevante bajo el plan gratuito utilizado: la URL pública generada cambia cada vez que el túnel se reinicia, lo que obliga a actualizar manualmente la URL configurada en la sección ``Aplicación de Prueba'' del panel de Partners de Tiendanube cada vez que ello ocurre. Esta fricción es asumida como parte del costo operativo de un entorno de desarrollo y no representa un problema para el prototipo, aunque sí sería un aspecto a resolver -mediante un dominio fijo o un hosting propio- de cara a una eventual publicación formal en la tienda de aplicaciones de Tiendanube, explícitamente señalada como trabajo futuro en el Alcance del proyecto.

\section{Arquitectura de red}
\label{sec:arquitectura-red}

La arquitectura de red del prototipo está determinada por el hecho de que la aplicación no se accede de forma directa, sino que se embebe como un \emph{iframe} dentro del panel de administración de Tiendanube, servido este último desde los servidores de la plataforma bajo HTTPS. Los navegadores modernos bloquean por política de \emph{contenido mixto} la carga de un \emph{iframe} servido por HTTP dentro de una página servida por HTTPS; en consecuencia, la aplicación embebida debe estar disponible también bajo una URL pública con HTTPS, incluso durante el desarrollo, donde de otro modo bastaría con \texttt{localhost}. Esta restricción es la que motiva el uso de ngrok, introducido en la sección anterior.

El flujo completo de comunicación, en el entorno de desarrollo, puede describirse en los siguientes pasos:

\begin{enumerate}
    \item Dentro del panel de administración de la tienda (servido por Tiendanube), el comerciante accede a la sección correspondiente a la aplicación, que Tiendanube renderiza como un \emph{iframe} cuyo origen (\texttt{src}) es la URL pública HTTPS provista por ngrok.

    \item ngrok reenvía el tráfico recibido en esa URL pública hacia el servidor de desarrollo de Vite, expuesto localmente en el puerto 5173, que sirve el frontend de React.

    \item Una vez cargado dentro del \emph{iframe}, el frontend inicializa el SDK Nexo, que establece comunicación con la ventana del panel de administración (la ventana padre) para obtener el contexto de la tienda y habilitar la sesión de la aplicación para ese comercio en particular.

    \item El frontend, aunque accedido a través de la URL pública de ngrok, se ejecuta en el navegador de la misma máquina donde corre el entorno de desarrollo. Por lo tanto, sus llamadas a la API propia de la aplicación se realizan directamente contra el backend en \texttt{localhost:3001}, sin necesidad de atravesar el túnel de ngrok.

    \item El backend, sobre la base del token obtenido mediante el flujo OAuth 2.0 descripto en la Sección~\ref{sec:estrategia-datos-ref} del Marco Teórico, consume los recursos correspondientes de la API REST de Tiendanube (\texttt{/orders}, \texttt{/products}, \texttt{/customers}, entre otros, detallados en esa misma sección) \parencite{TiendanubeAPI2024}, y devuelve al frontend los datos ya procesados.

    \item El frontend renderiza la respuesta dentro del \emph{iframe}, integrada visualmente al panel de administración mediante los componentes de Nimbus.
\end{enumerate}

Esta topología de red es propia del entorno de desarrollo y depende de una herramienta de terceros (ngrok) que introduce la fricción operativa descripta en la sección anterior. Reemplazar el túnel de desarrollo por un servicio de \emph{hosting} propio y permanente es, junto con la publicación formal en la tienda de aplicaciones de Tiendanube, uno de los aspectos que quedan explícitamente fuera del alcance del prototipo y planteados como trabajo futuro.

A la fecha de esta entrega, esta arquitectura fue validada de manera funcional de extremo a extremo: la aplicación carga correctamente embebida dentro del \emph{iframe} del panel de administración de la tienda de desarrollo de Tiendanube, con el intercambio (\emph{handshake}) de Nexo funcionando de forma correcta entre el \emph{iframe} y la ventana padre.

\section{Entorno de desarrollo}

Levantar el entorno de desarrollo completo requiere ejecutar tres procesos en simultáneo: el backend (\texttt{cd backend \&\& npm run dev}, puerto 3001), el frontend (\texttt{cd frontend \&\& npm run dev}, puerto 5173) y el túnel de ngrok (\texttt{ngrok http 5173}). Esta separación en tres procesos independientes refleja directamente la arquitectura de red descripta en la Sección~\ref{sec:arquitectura-red}: cada uno de ellos corresponde a un componente distinto del flujo de comunicación (servidor de API, servidor de la interfaz y túnel público).

La configuración de cada capa se administra mediante variables de entorno independientes en cada carpeta del monorepo. El backend concentra las credenciales sensibles de la aplicación (\texttt{Client ID} y \texttt{Client Secret} de Tiendanube), mientras que el frontend recibe únicamente el \texttt{Client ID}, dato público necesario para inicializar el SDK Nexo del lado del navegador. Esta separación no es incidental: mantiene el \texttt{Client Secret} fuera del código que se ejecuta en el navegador del usuario, en línea con el rol del backend como intermediario autenticado descripto en la Sección~\ref{sec:backend-node-express}. Cabe señalar, además, que Vite solo lee las variables de entorno del frontend al iniciar el proceso, por lo que una modificación en su archivo de configuración requiere reiniciar el servidor de desarrollo correspondiente para que los nuevos valores se apliquen.
